% ============================================================================
%  Grounded Regulatory Compliance Review: An Engineering Case Study
%
%  Build:  pdflatex compliance-engine.tex   (twice, for refs)
%  No external .bib required — references are inlined via thebibliography.
% ============================================================================
\documentclass[10pt,conference]{article}

\usepackage[utf8]{inputenc}
\usepackage[T1]{fontenc}
\usepackage[margin=1in]{geometry}
\usepackage{booktabs}
\usepackage{graphicx}
\usepackage{amsmath}
\usepackage{amssymb}
\usepackage{listings}
\usepackage{xcolor}
\usepackage{hyperref}
\usepackage{caption}
\hypersetup{colorlinks=true, linkcolor=black, citecolor=black, urlcolor=blue}

\lstset{
  basicstyle=\ttfamily\footnotesize,
  breaklines=true,
  frame=single,
  columns=fullflexible,
  showstringspaces=false,
  xleftmargin=0.5em,
}

\newcommand{\code}[1]{\texttt{\small #1}}

\title{\textbf{Grounded Regulatory Compliance Review:\\
An Engineering Case Study of a Precedent-Imitating\\
Retrieval System in Insurance Marketing}}

\author{%
  Bajaj Allianz Life Insurance --- Marketing Compliance\\
  \small \textit{Draft --- internal review}
}

\date{31 July 2026}

\begin{document}
\maketitle

\begin{abstract}
We describe the design, evolution and measured behaviour of a production
retrieval-augmented system that reviews life-insurance marketing copy against
IRDAI and SEBI obligations, brand policy, and product-specific factual
constraints. The system's central commitment is \emph{grounding before
generation}: no finding is emitted from unaided model judgement. Every
violation is anchored either in a retrieved past human-reviewer decision (a
``precedent''), an active rule, a deterministic product fact, a mandatory
disclosure obligation, or --- only above a confidence floor and with an
explicit regulatory basis --- a novel judgement. Anything the system cannot
substantively evaluate \emph{fails closed} to human review rather than being
recorded as compliant.

We report the system's architectural history across five months of
development, including several instances where a design was reversed after
production evidence contradicted it. We then present an honest measurement of
the deployed system over 29 documents and 1{,}537 findings, which surfaces
three material limitations: a corpus in which 58.4\,\% of precedents carry no
usable product scope, an aggregate scoring path that disagrees with its own
subscores, and the highest-firing rule family exhibiting zero reviewer-accepted
precision on a small sample. We argue that the reporting of such negative
results is the substantive contribution: the mechanisms that make this system
safe are inseparable from the mechanisms that make its failures visible.
\end{abstract}

\section{Introduction}

Marketing copy for regulated life-insurance products is reviewed by human
compliance officers before publication. The review is expensive, slow,
inconsistent across reviewers, and --- critically --- \emph{asymmetric in its
error costs}. A false positive wastes reviewer attention. A false negative
ships non-compliant claims about guaranteed returns, tax treatment or fund
performance to the public, with regulatory consequence.

This asymmetry is the design premise of the entire system. It explains choices
that would be unusual in a general-purpose retrieval system: relevance floors
applied per retrieval leg rather than post-fusion, findings that are persisted
but excluded from scoring rather than discarded, a deterministic (non-generative)
adjudicator for mandated disclosure wording, and a persistence gate that refuses
to record a grade for any run it cannot justify.

\subsection{Contributions}

\begin{enumerate}
  \item A \emph{grounding taxonomy} with five tiers of decreasing evidential
        strength, in which the tier is a first-class, queryable property of
        every finding rather than an implementation detail
        (Section~\ref{sec:grounding}).
  \item A \emph{fail-closed} persistence contract, enforced not once but at
        every upstream boundary, with the specific failure each guard prevents
        documented at the guard (Section~\ref{sec:failclosed}).
  \item A \emph{precedent memory} built from real reviewer decisions, with a
        canonicalisation scheme, two independent evaluation-leakage barriers,
        and a documented idempotency fragility arising from deriving a
        deduplication key from a generated field (Section~\ref{sec:precedent}).
  \item A \emph{bounded} feedback mechanism (Beta--Binomial rule reliability)
        together with a deliberate refusal to train on the evaluation metric
        (Section~\ref{sec:feedback}).
  \item A \emph{deterministic applicability layer} separating regulatory
        applicability from semantic similarity, motivated by a root-cause
        analysis of cross-product contamination (Section~\ref{sec:applicability}).
  \item An \emph{honest measurement} of the deployed system, including three
        negative results (Section~\ref{sec:evaluation}).
\end{enumerate}

\section{System Overview}

The pipeline is a six-node directed graph executed per submission:

\begin{center}
\code{preprocess} $\rightarrow$ \code{dispatch} $\rightarrow$ \code{analysis}
$\rightarrow$ \code{disclosure} $\rightarrow$ \code{scoring} $\rightarrow$ \code{persist}
\end{center}

\code{preprocess} extracts text, chunks it section-aware, and resolves product
identity (UIN, family, market segment). \code{dispatch} retrieves evidence per
chunk. \code{analysis} grades each chunk against retrieved evidence.
\code{disclosure} runs a separate deterministic mandatory-disclosure check.
\code{scoring} aggregates. Persistence is gated (Section~\ref{sec:failclosed}).

Retrieval is hybrid --- dense vector similarity and BM25 lexical search, fused
by Reciprocal Rank Fusion --- over six indexes: submission chunks, rules,
regulator source passages, v1 reviewer examples, v2 canonical precedents, and
approved product brochures.

\section{Grounding Before Generation}
\label{sec:grounding}

Each finding carries a \code{grounding} tier. The ranking is explicit in code:

\begin{lstlisting}
_TIER_RANK = {"disclosure": 5, "product_fact": 4,
              "precedent": 3, "rule": 2, "novel": 1}
\end{lstlisting}

\begin{table}[h]
\centering
\caption{Grounding tiers, strongest evidence first.}
\begin{tabular}{@{}llp{6.6cm}@{}}
\toprule
\textbf{Tier} & \textbf{Rank} & \textbf{Evidence} \\
\midrule
\code{disclosure}   & 5 & Verbatim match against a mandated disclaimer registry \\
\code{product\_fact} & 4 & Deterministic UIN-keyed product fact card \\
\code{precedent}    & 3 & A retrieved past human reviewer decision \\
\code{rule}         & 2 & An active, versioned rule \\
\code{novel}        & 1 & Model judgement, requiring an explicit regulatory basis and clearing a confidence floor \\
\bottomrule
\end{tabular}
\end{table}

The \code{novel} tier exists because of a measured coverage failure. In the
first precedent-only implementation, a chunk retrieving zero precedents
produced zero findings --- issues outside the corpus were missed entirely. The
design note recording this also documents a second, qualitative failure: the
system's output did not read like a reviewer. A representative bad output was
recorded verbatim in the design document as motivation:

\begin{quote}\itshape
``The new chunk mentions tax-free maturity benefits, which is similar to the
precedent that flagged a claim about life insurance.''
\end{quote}

This is a statement about retrieval, not about compliance. The reviewer-voice
redesign required each finding to name what is wrong, why it is wrong, and what
should replace it.

\section{Failing Closed}
\label{sec:failclosed}

The system may not record a grade for a document it did not substantively
evaluate. A run is blocked from persistence when: no chunks were produced;
graph status is failed; any degradation flag is set; or any chunk failed to
grade. Blocked runs are marked \code{needs\_review} (recoverable) or
\code{failed} (hard error) --- never scored.

The same posture is replicated at every boundary where a silent empty result
would be indistinguishable from a clean document. Three representative cases,
each documented at the guard:

\begin{itemize}
  \item An LLM returning an unparseable response must raise, because ``a
        fabricated empty result would be scored as a clean (100/A) document,
        i.e.\ the pipeline would fail OPEN.''
  \item A response truncated at the token ceiling
        (\code{finish\_reason='length'}) is rejected rather than partially
        accepted, since accepting it ``would silently under-report violations''
        and retrying cannot help.
  \item An empty rule set raises rather than returning \code{\{\}}, because an
        empty dict ``is indistinguishable from `zero rules configured','' and
        the pipeline would then grade with no rule grounding while appearing
        healthy.
\end{itemize}

\subsection{Suppression rather than deletion}

A design reversal worth recording. The original specification stated that novel
findings below the confidence floor were \emph{dropped before persistence}. The
shipped implementation inverts this: they are persisted with
\code{suppressed=True}, excluded from scoring, and surfaced in a separate human
review lane. The reasoning given in code is that a dropped low-confidence
finding is ``an undetectable false negative --- exactly the dangerous case in a
recall-critical compliance tool.''

Criticals are exempt from suppression entirely, on the grounds that ``false
negatives on criticals are the worst failure mode for a fail-closed compliance
tool.''

This is the asymmetry of Section~1 made operational: the system prefers a
visible, auditable false positive to an invisible false negative.

\section{Precedent Memory}
\label{sec:precedent}

\subsection{Origin}

The system's most consequential decision was made \emph{against} its brief. The
project was directed to replace a ``brittle rule-based RAG with no vector
infrastructure'' by building vector retrieval from scratch. An audit found the
premise false --- hybrid vector and BM25 retrieval with RRF fusion had been
present since the initial implementation. What did not exist was a corpus of
real reviewer decisions. The design note concludes:

\begin{quote}\itshape
``Therefore: we implement the brief's \emph{intent} using the \emph{existing}
architecture, not the brief's literal (duplicative, slightly-incompatible)
instructions.''
\end{quote}

The genuinely additive contribution was ingesting roughly 2{,}000--3{,}000
reviewed documents into tuples of the form \emph{(draft chunk, reviewer
comment, anchor text, final rewrite, reviewer, severity)}.

\subsection{Canonicalisation}

Version 2 of the corpus is issue-centric. Rather than embedding the raw text
pair, it embeds a \emph{signature} so that paraphrased violations still match:

\begin{lstlisting}
Issue:        {issue_type}
Why:          {why_rationale}
Flagged text: {highlighted_span}
Context:      {span_context}
\end{lstlisting}

Deduplication uses a canonical hash over the normalised triple
\emph{(issue\_type, highlighted\_span, reviewer\_comment)}, collapsing repeats
into one row with an \code{occurrence\_count}.

\subsection{An idempotency fragility}

The canonical hash is derived in part from \code{issue\_type}, which is
\emph{generated by a language model} during ingestion. Stable identifiers
across re-ingestion therefore depend on the enrichment cache persisting between
runs. The codebase states the consequence plainly: clearing the cache can
re-derive a different \code{issue\_type}, producing a new identifier and a
duplicate row. Cache persistence is documented as a hard operational
requirement.

We highlight this because it generalises: \emph{any} deduplication key derived
from generated output inherits that generator's nondeterminism. The
conventional assumption that a content hash is stable does not hold when the
content is itself sampled.

\subsection{Corpus hygiene and leakage}

Three hygiene layers operate: a one-time migration purging pure-response
comments (``done'', ``ok''), estimated at 12--15\,\% of a 5{,}323-row corpus; a
runtime thin-comment filter mirroring it, to protect against future bad
ingests; and a fallback to the v1 corpus only when v2 is empty.

Evaluation leakage is blocked twice, independently: retrieval excludes
precedents originating from the document under review, and the replay harness
splits train/eval deterministically by source-file hash. Neither barrier alone
is trusted.

\section{Bounded Feedback and the Refusal to Optimise the Metric}
\label{sec:feedback}

Each rule carries Beta pseudo-counts $(\alpha, \beta)$; its reliability is the
posterior mean
\[
\theta = \frac{\alpha}{\alpha + \beta},
\]
learned from reviewer accept/reject verdicts, and multiplied into the rule's
score deduction. Four safety properties are enforced:

\begin{enumerate}
  \item Absent feedback, $\theta = 1.0$ --- a rule with no evidence keeps full
        penalty. The system never silently under-penalises.
  \item A reliability floor (0.3) ensures feedback can discount a rule but
        never erase it: ``deactivating a rule is a human decision, not an
        emergent side effect of feedback.''
  \item The prior has strength 10, so a single verdict moves $\theta$ by at
        most $1/(\alpha+\beta+1)$. One reviewer click cannot whipsaw a rule.
  \item A changed verdict reverts the prior count before applying the new one.
\end{enumerate}

\subsection{The held-out reviewer score}

The system separately records a reviewer's own document-level score. This value
is \emph{never} an input to scoring or to weight updates. The stated reason,
repeated at four independent sites in the codebase, is that training on it
would Goodhart the metric that is supposed to demonstrate improvement. Its only
sanctioned use is the convergence curve
\[
\big|\, \text{score}_{\text{system}} - \text{score}_{\text{reviewer}} \,\big|
\]
observed over time.

We regard this as the system's most defensible design decision and its least
common one. The temptation to optimise directly against the available human
signal is strong, and yielding to it would have destroyed the only independent
evidence of improvement.

\subsection{What this mechanism is not}

The codebase is explicit that this is not continuous learning: there is no
approval gate, no training pipeline, and no model retraining. The corresponding
analytics surface returns an \code{insufficient\_data} shape rather than a
fabricated number when the underlying evidence does not exist. At the time of
measurement the reliability audit table existed but was never written to, so
the reliability-history panel correctly reported that it had nothing to show.

\section{Applicability: Separating Regulation from Similarity}
\label{sec:applicability}

\subsection{The failure}

A root-cause analysis found that product identity \emph{was} resolved before
retrieval and product scope metadata \emph{did} exist on the corpora --- but no
retrieval or post-retrieval step consumed it, in three separate places. Rule
retrieval filtered on regulator category (``IRDAI'') rather than product;
precedent retrieval passed no filter at all and additionally dropped the
product field before validation could occur; and a degraded-mode fallback
injected the entire active rule set into every chunk.

That last point is the sharpest: the fallback converted embedding outages into
\emph{maximum} cross-product contamination rather than reduced recall.

The analysis also isolates why relevance floors could not have caught this:
\begin{quote}\itshape
``cosine + ts\_rank floors exist but they are \emph{relevance} floors, not
\emph{applicability} floors --- an on-topic-but-wrong-product chunk sails over
them.''
\end{quote}

Semantic similarity cannot express regulatory applicability. A ULIP
risk-disclosure precedent and a term-plan risk-disclosure precedent are
near-neighbours in embedding space and are not interchangeable in law.

\subsection{The contract}

Applicability is decided deterministically, before similarity is allowed to
matter, under an explicit contract: a conflicting product category is rejected;
untagged or unmappable items are treated as global and accepted; an unresolved
product scope rejects nothing; a unit-linked rider widens its scope; and every
decision --- accept or reject, with reason --- is recorded per candidate.

Crucially the default is \emph{permissive}. A heuristic ingest tag must never
over-block retrieval, because over-blocking silently removes evidence and
manifests as a false negative.

\subsection{Cross-cutting categories}

Not every tag denotes a product family. In this corpus \code{child} is a
variant of a unit-linked product \emph{and} a rider benefit attached to term
plans; scoping it to either family would hide it from the other. Such tags are
now named explicitly as cross-cutting rather than reaching the ``global''
outcome by accident through the unmappable-tag path --- a distinction with no
behavioural difference today, but one that prevents a later vocabulary fix from
silently narrowing retrieval.

\subsection{Why metadata, not embedding}

Product scope is stored beside the vector, not inside the embedded text. This
is deliberate. Embedding a product name into every precedent signature would
cause precedents to cluster by product rather than by \emph{issue type}, which
is what retrieval must discriminate. Categorical scoping is an exact-match
problem and belongs in a filter; embedding is for semantic similarity. A
practical corollary is that re-tagging a corpus costs nothing in re-embedding.

\section{Deterministic Adjudication of Mandated Wording}
\label{sec:disclosure}

Mandatory disclosures are handled by a separate path: a registry of approved
wording, a trigger layer (deterministic product-line and keyword rules, unioned
with an optional model backstop for paraphrased obligations), and a
\emph{deterministic} matcher.

The division of labour is a legal-defensibility argument rather than a quality
one. Whether a document \emph{incurs} an obligation is a recall problem where
paraphrase understanding helps, and a model is used. Whether the document's
wording \emph{satisfies} the obligation is a verbatim-match question with a
known correct answer, and no generative model adjudicates it. The model may
only add obligations, never remove a deterministically-established one.

\subsection{A diagnostic case study}

A creative visibly containing an approved sentence was reported as ``present
but altered (similarity 0.50).'' The investigation concluded the matcher was
not at fault: the sentence never reached it, having been dropped during text
extraction, after which every downstream stage behaved to specification and
compounded the error into a misleading message.

Two diagnostic steps are worth reproducing. First, the finding's provenance was
recorded as originating from the model backstop, which is assigned only when
the deterministic trigger did \emph{not} fire --- proving the regular
expression never saw the phrase. Second, and more striking, the similarity
score itself was identified as a fingerprint of \emph{absence}: unrelated
documents containing no disclaimer at all score in the 0.48--0.55 band, because
that is the floor produced by incidental overlap of common English tokens. A
score of 0.50 therefore indicated ``not present'' rather than ``present but
altered.''

The root cause was that the DOCX extractor iterated only document paragraphs,
which excludes headers, footers, text boxes and table cells --- precisely where
marketing creatives place mandated disclaimers.

The remediation included a consequence beyond the reported symptom: the matcher
now protects legally-significant tokens (\emph{not}, \emph{no}, \emph{may},
\emph{guaranteed}, \emph{risk}, \emph{past}) from being treated as ignorable.
Previously, a disclaimer whose meaning had been \emph{inverted} by deleting
``not'' scored 0.96 and passed as present --- undetected tampering.

\section{Measurement}
\label{sec:evaluation}

\subsection{Baseline}

The only quantitative baseline recorded during development, over 9--10
documents, is reproduced below. Its own design note names
\code{precision\_presence} as the canary: a regression there would indicate
novel findings hallucinating.

\begin{table}[h]
\centering
\caption{Development baseline, 27 May 2026.}
\begin{tabular}{@{}lr@{}}
\toprule
\textbf{Metric} & \textbf{Value} \\
\midrule
precision\_presence & 0.21 \\
recall\_presence    & 0.66 \\
anchor\_recall      & 0.35 \\
missed\_criticals   & 32 \\
generated / real    & 3.0$\times$ \\
\bottomrule
\end{tabular}
\end{table}

The recall-over-precision profile is intentional given the error asymmetry, but
a $3\times$ generation ratio is a real reviewer-attention cost.

\subsection{Deployed system}

Measured on the live deployment, 31 July 2026: 29 submissions, 1{,}537
findings, 42 completed analyses, 2{,}440 precedents, 148 active rules.

\begin{table}[h]
\centering
\caption{Deployed characteristics.}
\begin{tabular}{@{}lr@{}}
\toprule
\textbf{Property} & \textbf{Value} \\
\midrule
Mean analysis latency & 133{,}910\,ms \\
Median (p50) & 128{,}520\,ms \\
p95 & 268{,}293\,ms \\
Findings of Medium severity & 80\,\% \\
Grade distribution (F / D / C / B / A) & 36 / 2 / 1 / 3 / 0 \\
Reviewer verdicts recorded & 22 of 1{,}308 (1.7\,\%) \\
\bottomrule
\end{tabular}
\end{table}

\subsection{Negative result 1: scope coverage}

Of 2{,}440 precedents, only 41.6\,\% carry a tag that resolves to a product
family:

\begin{table}[h]
\centering
\caption{Precedent scope coverage.}
\begin{tabular}{@{}lrrl@{}}
\toprule
\textbf{Tag} & \textbf{Count} & \textbf{\%} & \textbf{Effect} \\
\midrule
(untagged) & 1{,}167 & 47.8 & global \\
Term       & 429 & 17.6 & scoped \\
ULIP       & 333 & 13.6 & scoped \\
Child      & 259 & 10.6 & global (cross-cutting) \\
Pension    & 178 &  7.3 & scoped \\
Savings    &  74 &  3.0 & scoped \\
\midrule
\textbf{Effectively global} & \textbf{1{,}426} & \textbf{58.4} & \\
\bottomrule
\end{tabular}
\end{table}

Two observations. First, the applicability layer is correct but
under-supplied: the majority of the corpus is exempt from the scoping it
implements. Second, the corpus contains \emph{no} instances of the
participating/non-participating distinction, because those ingest heuristics
postdate the corpus. The market segmentation the business considers primary is
therefore unexpressed at the precedent level.

This also bounds an apparently attractive optimisation. Pushing the scope
filter into SQL --- retrieving within scope rather than filtering afterwards
--- yields less than expected when 58.4\,\% of rows must be retained
regardless.

\subsection{Negative result 2: aggregate--subscore disagreement}

The deployed report surface displays an overall score of 0.0 alongside
component scores of 35.8--99.1 for the same document. No weighting of those
components produces 0.0. Either the aggregate is not derived from the
components, or an undisclosed override zeroes it. The population is consistent
with a systematic effect: 36 of 42 analyses graded F, mean score 14.2.

The same document additionally reports different finding counts on different
surfaces (53 versus 68), most plausibly the suppressed review lane being
included in one query and excluded in another --- a direct, and predictable,
cost of the deliberate decision in Section~\ref{sec:failclosed} to persist
suppressed findings rather than discard them. Introducing a second class of
finding requires every consumer to state which class it counts.

\subsection{Negative result 3: rule precision}

The highest-firing rules --- brand-name conventions, 45, 23 and 9 activations
--- have been rejected by reviewers in every instance so far adjudicated (0\,\%
precision on 4 and 2 verdicts respectively). With 22 verdicts against 1{,}308
findings, this is a signal to investigate rather than a conclusion. But the
direction is consistent with the $3\times$ over-generation of the development
baseline and with Medium severity comprising 80\,\% of all findings: the system
is producing high-volume, low-specificity output in one rule family.

Notably, this was surfaced \emph{by the system's own analytics}, because
reviewer verdicts are recorded per rule with provenance. The instrumentation
built to demonstrate improvement was the instrumentation that revealed the
defect.

\section{Discussion}

\subsection{Documentation drift as a failure mode}

An audit of the codebase against its own documentation found 15 discrepancies,
including a pipeline described as five nodes when it has six, four different
active-rule counts across seven interface surfaces, a knowledge-base statistic
that reports a rendering cap as a corpus size, and three load-bearing code
comments citing design documents that never existed in version control.

We include this because it is a genuine reliability concern rather than an
editorial one. In a regulated setting, ``which model graded this document'' is
an audit question, and the deployed system answered it two different ways on
two different screens.

\subsection{On negative results}

The mechanisms that make this system safe --- per-candidate accept/reject
provenance, held-out evaluation, honest \code{insufficient\_data} responses,
suppressed-but-persisted findings --- are the same mechanisms that made the
three negative results in Section~\ref{sec:evaluation} discoverable. A system
built to look successful would have reported a mean compliance score and
stopped.

\section{Limitations}

The evaluation is observational, on a single deployment, with 1.7\,\% reviewer
adjudication coverage; precision figures are indicative only. No inter-rater
agreement was measured. No ablation isolates the contribution of individual
grounding tiers. The development baseline covers 9--10 documents. Cost
telemetry was not functioning during measurement, so efficiency claims are not
made. The system has not been compared against a human-reviewer control or a
non-grounded baseline; such a comparison is the obvious next work.

\section{Conclusion}

Grounding before generation, fail-closed persistence, deterministic
adjudication of mandated wording, bounded feedback, and a refusal to train on
the evaluation metric together produce a compliance reviewer whose findings are
auditable and whose failures are visible. Measurement of the deployed system
shows the architecture is sound and the \emph{supply} to it is not: an
applicability layer starved of tags, an aggregate score disagreeing with its
components, and a rule family generating findings reviewers consistently
reject. These are tractable problems, and they were found because the system
was instrumented to be falsifiable.

\begin{thebibliography}{9}
\footnotesize

\bibitem{rrf}
G.~V. Cormack, C.~L.~A. Clarke, and S.~B\"uttcher.
\newblock Reciprocal rank fusion outperforms Condorcet and individual rank
  learning methods.
\newblock In \emph{SIGIR}, 2009.

\bibitem{bm25}
S.~Robertson and H.~Zaragoza.
\newblock The probabilistic relevance framework: BM25 and beyond.
\newblock \emph{Foundations and Trends in Information Retrieval}, 3(4), 2009.

\bibitem{rag}
P.~Lewis et al.
\newblock Retrieval-augmented generation for knowledge-intensive NLP tasks.
\newblock In \emph{NeurIPS}, 2020.

\bibitem{goodhart}
M.~Strathern.
\newblock `Improving ratings': audit in the British University system.
\newblock \emph{European Review}, 5(3), 1997.

\bibitem{irdai}
Insurance Regulatory and Development Authority of India.
\newblock Guidelines on advertisements, promotion and publicity of insurance
  companies and insurance intermediaries.

\end{thebibliography}

\end{document}
